\documentclass[11pt,a4paper]{article}

% ---------- Packages ----------
\usepackage[utf8]{inputenc}
\usepackage[T1]{fontenc}
\usepackage[margin=2.5cm]{geometry}
\usepackage{mathpazo}
\usepackage{microtype}
\usepackage{amsmath,amssymb}
\usepackage{booktabs}
\usepackage{array}
\usepackage{enumitem}
\usepackage{xcolor}
\usepackage{titlesec}
\usepackage{setspace}
\usepackage[hidelinks]{hyperref}

\setstretch{1.08}
\setlength{\parskip}{0.45em}
\setlength{\parindent}{0pt}
\titleformat{\section}{\large\bfseries}{\thesection}{0.7em}{}
\titleformat{\subsection}{\normalsize\bfseries}{\thesubsection}{0.7em}{}
\definecolor{note}{RGB}{90,90,90}
\newcommand{\hits}{\textsc{Hits}@100}
\newcommand{\cn}{\mathrm{CN}}
\renewcommand{\arraystretch}{1.2}

% ---------- Title ----------
\title{\vspace{-1em}\textbf{Common-Neighbour Completion for Link Prediction:\\
A Literature Review Centred on IECNC}}
\author{%
  \begin{tabular}{@{}ccc@{}}
    \textbf{Shruten Pande} & \textbf{Hitansh Agarwal} & \textbf{Sumantu Sharma}\\
    230957052 & 230957218 & 230957100
  \end{tabular}}
\date{}

\begin{document}
\maketitle
\vspace{-1em}

\begin{abstract}
\noindent
Link prediction asks whether an unobserved edge exists between two nodes of a network.
This review uses \emph{Common Neighbor Completion with Information Entropy} (IECNC)
\cite{zhou2025iecnc} as its anchor and situates it within the line of work that moves from
hand-crafted proximity heuristics, through learned subgraph structure (SEAL)
\cite{zhang2018seal} and sketch-based efficiency (ELPH/BUDDY) \cite{chamberlain2023elph},
to neural common-neighbour pooling and completion (NCN/NCNC) \cite{wang2024ncnc}.
We summarise each method, tabulate the evidence IECNC reports against its baselines,
and examine the contribution of its entropy term critically. We also note
a bibliometric study of hybrid text--citation clustering \cite{janssens2007} as a
thematically adjacent, but methodologically distinct, work. We close with open issues,
including internal inconsistencies in the IECNC text that readers should check against its tables.
\end{abstract}

% =====================================================================
\section{Introduction}
Link prediction underpins friend recommendation, knowledge-graph completion, biological
interaction discovery and citation analysis. Classical approaches score a candidate pair
$(i,j)$ with a hand-designed function of the graph---most prominently the number of common
neighbours $\cn(i,j)=|N(i)\cap N(j)|$ and its weighted variants. Graph neural networks (GNNs)
promised to replace such hand-crafting, yet message-passing GNNs often under-perform simple
heuristics on link prediction because node-level embeddings cannot count triangles
or separate structurally equivalent nodes \cite{chamberlain2023elph}.
A recent research thread therefore injects \emph{pairwise structural information}---especially
common neighbours---into neural models while keeping them scalable.

IECNC \cite{zhou2025iecnc} sits at the end of this thread. It builds directly on NCNC
\cite{wang2024ncnc}, retaining the idea of completing a possibly incomplete common-neighbour
set, and adds (i)~explicit logical neighbourhood sets (intersection, union, directional
differences) and (ii)~an information-entropy weighting of the predicted common-neighbour
probabilities. The remainder of this review follows that lineage and ends with a critical
reading of IECNC's evidence.

% =====================================================================
\section{Background: From Heuristics to Learned Structure}
\subsection{Heuristics and their limits}
Heuristic scores such as Common Neighbours, Adamic--Adar (AA) and Resource Allocation (RA)
are simple, interpretable and scalable, but each encodes a fixed assumption about when two
nodes should link. SEAL's authors argue that this limits them on networks where the
assumption fails, and propose learning a suitable heuristic from the data instead
\cite{zhang2018seal}. In IECNC's own evaluation the three heuristics are by far the weakest
baselines (e.g.\ \hits{} of 33.92, 39.85 and 41.07 for CN, AA and RA on Cora, versus 66.79
for GCN) \cite{zhou2025iecnc}.

\subsection{Message passing and its expressiveness gap}
An MPNN updates each node from its neighbours' messages over several layers, and a link is
scored from the two endpoint representations \cite{gilmer2017mpnn}. IECNC adopts this view,
writing the link score as a function of $h_i$, $h_j$ and a pairwise context $C_{ij}$
\cite{zhou2025iecnc}. The expressiveness limits of plain MPNN scoring are the motivation
shared by SEAL, ELPH/BUDDY and NCN: all add pairwise or subgraph-level structure.

% =====================================================================
\section{The Lineage Leading to IECNC}
\subsection{SEAL: learning heuristics from enclosing subgraphs}
SEAL extracts the local subgraph around each target link, labels its nodes with
double-radius node labelling, and classifies the subgraph with a graph-level GNN; node
embeddings and attributes may optionally be concatenated \cite{zhang2018seal}. Its theoretical
contribution is the $\gamma$-decaying heuristic theory, which argues that a wide family of
high-order heuristics can be well approximated from small local subgraphs. SEAL was evaluated on eight
networks (USAir, NS, PB, Yeast, C.ele, Power, Router, E.coli) with AUC and average precision, and the authors
report that it generally outperformed heuristic, latent-feature and embedding baselines,
including WLNM. The cost is that a subgraph must be built and processed for every candidate link.

\subsection{ELPH and BUDDY: sketching instead of subgraphs}
Chamberlain et al.\ analyse the components of subgraph GNNs and propose ELPH, a full-graph GNN that passes
HyperLogLog and MinHash sketches as messages, approximating the key structural quantities
of subgraph methods without constructing subgraphs, and is provably more expressive than
MPNNs \cite{chamberlain2023elph}. BUDDY precomputes the sketches and diffused features so
that prediction reduces to an MLP. The authors report BUDDY to be 200--1000$\times$
faster than SEAL in dynamic mode on the citation dataset, while remaining competitive
or better on Planetoid and OGB benchmarks. Their contribution is therefore primarily one of \emph{efficiency at
comparable accuracy}.

\subsection{NCN and NCNC: pooling and completing common neighbours}
Wang et al.\ classify prior models as SF-then-MPNN (structural features augment the MPNN
input) or SF-and-MPNN (decoupled), and propose \emph{MPNN-then-SF}: run the MPNN once on the
input graph and use the structural common-neighbour set to \emph{guide the pooling} of
node representations for each candidate pair. This is Neural Common Neighbor (NCN)
\cite{wang2024ncnc}. They then study \emph{graph incompleteness}: since the prediction targets are
by definition unobserved, the input graph lacks edges, which reduces the number of common neighbours and
induces a train/test distribution shift. Neural Common Neighbor with Completion (NCNC)
uses a link-prediction model's predicted edge probabilities to complete the common-neighbour
structure before a second, NCN-style prediction. Evaluation covers seven benchmarks
(Cora, Citeseer, PubMed, ogbl-collab, ogbl-ppa, ogbl-citation2, ogbl-ddi), with 10 random seeds,
ablations, and runtime/memory comparisons. IECNC adopts NCNC's completion principle
as its starting point.

\begin{table}[t]
\centering
\caption{Summary of the surveyed link-prediction methods.}
\label{tab:methods}
\small
\begin{tabular}{@{}p{2.1cm}p{4.9cm}p{3.7cm}p{3.6cm}@{}}
\toprule
\textbf{Method} & \textbf{Core mechanism} & \textbf{Structural signal} & \textbf{Main claim}\\
\midrule
SEAL \cite{zhang2018seal} & Enclosing subgraph, double-radius labelling, graph-level GNN & Learned from local subgraph ($\gamma$-decay theory) & Learned heuristics beat fixed ones\\
ELPH / BUDDY \cite{chamberlain2023elph} & HyperLogLog/MinHash sketches as messages; precomputed in BUDDY & Approximate subgraph counts & Subgraph-level accuracy at far lower cost\\
NCN \cite{wang2024ncnc} & MPNN run once; pool over common neighbours (MPNN-then-SF) & Observed common neighbours & Expressive and scalable\\
NCNC \cite{wang2024ncnc} & NCN plus completion by predicted edge probabilities & Observed + predicted common neighbours & Mitigates graph incompleteness\\
IECNC \cite{zhou2025iecnc} & Logical neighbour sets, completion probabilities, entropy weighting & Completed common neighbours, uncertainty-weighted & Further gains over NCNC\\
\bottomrule
\end{tabular}
\end{table}

% =====================================================================
\section{IECNC in Detail}
\subsection{Method}
IECNC \cite{zhou2025iecnc} models a node pair $(i,j)$ through a \emph{logical context}
built from neighbour sets $q_i=N(i)$ and $q_j=N(j)$:
\begin{equation}
C_{ij}=\{\,q_i,\; q_j,\; q_i\wedge q_j,\; q_i\vee q_j,\; q_i-q_j,\; q_j-q_i\,\},
\end{equation}
where $\wedge$, $\vee$ and $-$ denote intersection, union and difference. Following the
paper's notation, each set is encoded by an MLP over the MPNN features of its members, and
the concatenation $[h_i;h_j;Q_\vee;Q_\wedge;Q_{i-j};Q_{j-i}]$ forms $C_{ij}$.

Completion is implemented through a per-node probability $P_{uij}$ that $u$ acts as a
common neighbour of $i$ and $j$:
\begin{equation}
P_{uij}=
\begin{cases}
1, & u\in A_\cap,\\
\operatorname{softmax}(Q_{i-j}), & u\in A_{i-j},\\
\operatorname{softmax}(Q_{j-i}), & u\in A_{j-i},\\
0, & \text{otherwise,}
\end{cases}
\end{equation}
so that observed common neighbours receive probability~1, nodes adjacent to only one
endpoint receive a learned probability, and the rest receive~0. The earlier, entropy-free
variant (PPNC) aggregates $\sum_u P_{uij}\,\mathrm{MPNN}(u,M,C_{ij})$. IECNC's distinctive step
replaces the probability weight with an information-entropy weight:
\begin{equation}
\mathrm{IECNC}(i,j)=-\sum_{u\in q}P_{uij}\ln P_{uij}\cdot \mathrm{MPNN}(u,M,C_{ij}).
\end{equation}
The authors motivate this as quantifying the uncertainty of each node's membership in the common-neighbour
set. They also state a design simplification: because higher-order information appeared to be
implicitly captured by first-order neighbours, only first-order neighbours and intersection-based
set aggregation are used.

\subsection{Experimental evidence}
IECNC is evaluated on five graphs: three Planetoid citation networks (Cora, Citeseer, PubMed)
and two Twitch gamer networks (Twitch\_RU, Twitch\_PT), using \hits{} and AUC, against
CN, AA, RA, GCN, SAGE, SEAL, NBFNet, Neo-GNN, BUDDY, NCN and NCNC \cite{zhou2025iecnc}.
Table~\ref{tab:results} reproduces the headline \hits{} comparison against the strongest baselines
from the paper's Table~1.

\begin{table}[t]
\centering
\caption{\hits{} (\%, mean $\pm$ s.d.) from Table~1 of \cite{zhou2025iecnc}; best per column in bold.
The last row is our arithmetic on the reported means (percentage points).}
\label{tab:results}
\small
\begin{tabular}{@{}lccccc@{}}
\toprule
\textbf{Method} & \textbf{Cora} & \textbf{Citeseer} & \textbf{PubMed} & \textbf{Twitch\_RU} & \textbf{Twitch\_PT}\\
\midrule
SEAL    & 81.71 $\pm$ 1.30 & 83.89 $\pm$ 2.15 & 75.54 $\pm$ 1.32 & 50.74 $\pm$ 0.43 & 46.80 $\pm$ 3.16\\
Neo-GNN & 80.42 $\pm$ 1.31 & 84.67 $\pm$ 2.16 & 73.93 $\pm$ 1.19 & 50.52 $\pm$ 0.37 & 44.13 $\pm$ 0.60\\
BUDDY   & 88.00 $\pm$ 0.44 & 92.93 $\pm$ 0.27 & 74.10 $\pm$ 0.78 & 55.94 $\pm$ 0.58 & 47.85 $\pm$ 0.20\\
NCN     & 89.05 $\pm$ 0.96 & 91.56 $\pm$ 1.43 & 79.05 $\pm$ 1.16 & 57.76 $\pm$ 0.87 & 51.19 $\pm$ 0.85\\
NCNC    & 89.65 $\pm$ 1.36 & 93.47 $\pm$ 0.95 & 81.29 $\pm$ 0.95 & 58.61 $\pm$ 0.71 & 52.42 $\pm$ 0.73\\
\midrule
IECNC   & \textbf{94.68 $\pm$ 1.43} & \textbf{96.44 $\pm$ 1.63} & \textbf{87.56 $\pm$ 1.51} & \textbf{62.18 $\pm$ 2.4} & \textbf{56.27 $\pm$ 1.34}\\
$\Delta$ vs.\ NCNC & +5.03 & +2.97 & +6.27 & +3.57 & +3.85\\
\bottomrule
\end{tabular}
\end{table}

For AUC, the paper reports IECNC at 97.70\% (Cora), 99.23\% (Citeseer) and 95.26\% (Twitch\_PT),
versus NCNC at 95.70\%, 97.66\% and 94.67\%; these values appear in the text and Fig.~3 only.
The gain over NCNC is thus consistently positive, and largest on PubMed.

\subsection{Ablations and the role of entropy}
The paper's ablations (its Tables~2--3) vary the activation function (sigmoid vs.\ softmax),
the completion strategy (probabilistic completion, PC, vs.\ information-entropy completion, IEC) and the components of $P_{uij}$.
Removing the entire completion step ($P_{uij}$) causes the largest drop, followed by removing
$A_\cap$, $A_{i-j}$ or $A_{j-i}$. Table~\ref{tab:entropy} isolates the entropy term from those ablations.

\begin{table}[t]
\centering
\caption{Effect of replacing probabilistic completion (PC) by entropy completion (IEC),
\hits{} (\%), from Table~2 of \cite{zhou2025iecnc}; $\Delta$ is computed from the reported means.}
\label{tab:entropy}
\small
\begin{tabular}{@{}llccccc@{}}
\toprule
\textbf{Activation} & & \textbf{Cora} & \textbf{Citeseer} & \textbf{PubMed} & \textbf{Twitch\_RU} & \textbf{Twitch\_PT}\\
\midrule
softmax & PC  & 93.78 & 94.43 & 82.54 & 58.64 & 54.41\\
softmax & IEC & 94.68 & 96.44 & 87.56 & 62.18 & 56.27\\
\multicolumn{2}{@{}l}{\emph{$\Delta$ (softmax)}} & +0.90 & +2.01 & +5.02 & +3.54 & +1.86\\
\midrule
sigmoid & PC  & 89.36 & 93.99 & 81.20 & 59.04 & 54.34\\
sigmoid & IEC & 89.96 & 93.97 & 81.69 & 57.07 & 53.71\\
\multicolumn{2}{@{}l}{\emph{$\Delta$ (sigmoid)}} & +0.60 & $-$0.02 & +0.49 & $-$1.97 & $-$0.63\\
\bottomrule
\end{tabular}
\end{table}

Two observations follow. First, with softmax the entropy term helps on every dataset, most
visibly on PubMed. Second, with sigmoid the benefit is negligible or negative on three of five
datasets, so the contribution of entropy appears to depend on the softmax parameterisation
rather than being an independent gain. The authors themselves note that with sigmoid the two completion
strategies differ little.

% =====================================================================
\section{Critical Reading of IECNC}
The following points are offered as reading notes, not as a replacement for checking the original paper.

\paragraph{Inconsistent headline numbers.}
The prose of Section~5.2 reports \hits{} of 94.09\% (Cora), 95.34\% (Citeseer), 85.72\% (PubMed) and a Twitch\_RU
standard deviation of 4.4, whereas Tables~1--3 and the Conclusion report 94.68\%, 96.44\%, 87.56\% and 2.4.
The three tables are mutually consistent, so this review uses the table values; the discrepancy should
nevertheless be resolved against the authors' code or results. The Fig.~4 captions also alternate between \textsc{Hits}@10 and \textsc{Hits}@3 for the lower panels.

\paragraph{What the entropy weight does.}
The paper motivates $-P\ln P$ as an uncertainty measure that is maximal near $P=0.5$. Two details
deserve scrutiny. (i)~Eq.~(2) assigns observed common neighbours $P=1$, for which $-P\ln P=0$; as written, those
nodes would contribute nothing to the score, yet the $A_\cap$ ablation shows they matter, suggesting that the
implementation differs from the equation or that the role of $A_\cap$ is mediated elsewhere. (ii)~Softmax
over a neighbour set yields small probabilities (Table~4 of the paper lists $P\approx10^{-4}$--$10^{-3}$); in that regime
$-P\ln P$ is a smooth, monotonically increasing function of $P$, so it behaves as a re-weighting of the
probability rather than a distinct uncertainty signal. The paper's Table~4 shows that near-identical probabilities map to
slightly different entropy values, which demonstrates a transformation, not that the transformation improves accuracy;
the ablation in Table~\ref{tab:entropy} is the stronger evidence.

\paragraph{Scope of evaluation.}
Despite the ``social networks'' framing, three of five datasets are citation graphs; only the two Twitch graphs are social.
No large OGB benchmarks (ogbl-collab, -ppa, -citation2, -ddi), which NCNC and BUDDY use, are reported, and
the paper claims scalability and efficiency without presenting runtime or memory measurements,
unlike NCNC \cite{wang2024ncnc} and BUDDY \cite{chamberlain2023elph}. The Cora--PubMed figures for BUDDY, NCN and NCNC in its Table~1 coincide with those
in the NCNC paper, indicating that baseline numbers were taken from prior work rather than re-run,
which is acceptable practice but means protocol parity should be confirmed.

\paragraph{Metric heterogeneity across the literature.}
SEAL reports AUC and AP on random 10\% holdouts of small networks \cite{zhang2018seal}; BUDDY and NCNC emphasise ranking
metrics such as \textsc{Hits}@$K$ on Planetoid and OGB splits \cite{chamberlain2023elph,wang2024ncnc}. Scores are therefore comparable only within a
given protocol, which limits cross-paper conclusions.

% =====================================================================
\section{A Thematically Adjacent Study: Hybrid Text--Citation Clustering}
Janssens, Gl\"anzel and De~Moor \cite{janssens2007} analyse 7{,}401 bioinformatics publications
(Web of Science and MEDLINE, 1981--2004) across seven time windows. They combine text similarity
(TF--IDF/LSI) with bibliographic-coupling distances through Fisher's inverse chi-square, cluster with Ward's
method, and select the number of clusters by combining dendrogram inspection, silhouette curves and
subsample stability. The result is nine static clusters linked into eleven evolving cluster chains.
This is a descriptive bibliometric analysis and \emph{not} a link-prediction method, and
it offers no comparison with the link-prediction baselines above. Its relevance is conceptual: it
supports the broader view that combining content with graph structure is more informative than either
alone, the same intuition behind attribute-aware structural link prediction. IECNC lists bioinformatics
among its application domains \cite{zhou2025iecnc}, but the two works share no method or benchmark.

% =====================================================================
\section{Open Issues and Directions}
\begin{itemize}[leftmargin=1.5em,itemsep=0.2em]
  \item \textbf{Isolating the entropy term.} A controlled comparison against other monotone reweightings of $P_{uij}$ (and against the sigmoid setting)
        would show whether entropy specifically, rather than extra weighting capacity, drives the gains.
  \item \textbf{Scalability evidence.} Runtime and memory comparisons with NCNC and BUDDY on OGB-scale graphs are needed to support efficiency claims.
  \item \textbf{Broader and temporal benchmarks.} Evaluation on genuinely social, evolving graphs would test the ``social network'' framing.
  \item \textbf{Completion quality.} NCNC and IECNC both rely on a first-stage predictor to complete the graph; error propagation from that stage is not analysed in the text of either paper.
  \item \textbf{Reproducibility.} Resolving the number discrepancies noted in Section~5 and releasing code would settle the headline results.
\end{itemize}

% =====================================================================
\section{Conclusion}
The trajectory from heuristics to SEAL, ELPH/BUDDY, NCN/NCNC and IECNC reflects a steady effort to give neural link predictors
the pairwise structural information that simple heuristics exploit, at decreasing computational cost
and with growing attention to missing edges. IECNC extends NCNC with logical neighbour sets and an entropy
weighting and reports consistent \hits{} gains of roughly 3--6 points over NCNC on five datasets. The ablations support the value of
completion and, with softmax, of the entropy term, but the evidence base is narrow, the entropy formulation raises
questions, and the text contains inconsistencies that should be resolved before the figures are cited as definitive.

% =====================================================================
\begin{thebibliography}{99}\small
\setlength{\itemsep}{0.2em}

\bibitem{zhou2025iecnc}
Z.~Zhou, G.~Wan, and B.~Du.
\newblock Common neighbor completion with information entropy for link prediction in social networks.
\newblock \emph{Data Science and Engineering}, 10:40--53, 2025.
\newblock doi:10.1007/s41019-024-00267-6.

\bibitem{wang2024ncnc}
X.~Wang, H.~Yang, and M.~Zhang.
\newblock Neural common neighbor with completion for link prediction.
\newblock In \emph{International Conference on Learning Representations (ICLR)}, 2024.

\bibitem{chamberlain2023elph}
B.~P. Chamberlain, S.~Shirobokov, E.~Rossi, F.~Frasca, T.~Markovich, N.~Hammerla, M.~M. Bronstein, and M.~Hansmire.
\newblock Graph neural networks for link prediction with subgraph sketching.
\newblock In \emph{International Conference on Learning Representations (ICLR)}, 2023.
\newblock arXiv:2209.15486.

\bibitem{zhang2018seal}
M.~Zhang and Y.~Chen.
\newblock Link prediction based on graph neural networks.
\newblock In \emph{Advances in Neural Information Processing Systems (NeurIPS)}, 31, 2018.

\bibitem{janssens2007}
F.~Janssens, W.~Gl\"anzel, and B.~De~Moor.
\newblock Dynamic hybrid clustering of bioinformatics by incorporating text mining and citation analysis.
\newblock In \emph{Proc.\ ACM SIGKDD International Conference on Knowledge Discovery and Data Mining (KDD'07)}, San Jose, CA, 2007.

\bibitem{gilmer2017mpnn}
J.~Gilmer, S.~S. Schoenholz, P.~F. Riley, O.~Vinyals, and G.~E. Dahl.
\newblock Neural message passing for quantum chemistry.
\newblock In \emph{International Conference on Machine Learning (ICML)}, pages 1263--1272, 2017.

\end{thebibliography}

\end{document}
